% ═══════════════════════════════════════════════════════════════
% LERNBLATT - DS-02: Mi nueva casa – Repaso
% Klasse 9, Unidad 2
% ═══════════════════════════════════════════════════════════════

\documentclass[11pt, a4paper]{article}

\usepackage[utf8]{inputenc}
\usepackage[T1]{fontenc}
\usepackage[ngerman]{babel}

% --- SCHRIFTEN ---
\usepackage{helvet}
\renewcommand{\familydefault}{\sfdefault}
\usepackage{palatino}
\newcommand{\seriftext}[1]{{\fontfamily{ppl}\selectfont #1}}

\usepackage{microtype}
\usepackage[margin=1.8cm, top=2cm, bottom=1.5cm]{geometry}
\usepackage{enumitem}
\usepackage{tabularx}
\usepackage{booktabs}
\usepackage{array}
\usepackage{multicol}
\usepackage{tikz}

\usepackage{fancyhdr}
\usepackage{lastpage}

% ═══════════════════════════════════════════════════════════════
% VARIABLEN
% ═══════════════════════════════════════════════════════════════

\newcommand{\lektion}{02}
\newcommand{\thema}{Mi nueva casa -- Repaso}
\newcommand{\untertitel}{Wiederholung + Pretérito perfecto}
\newcommand{\klasse}{9}
\newcommand{\lernziele}{Vokabelwiederholung · Bewegungsverben · Pretérito perfecto}

% ═══════════════════════════════════════════════════════════════
% FARBEN
% ═══════════════════════════════════════════════════════════════

\usepackage{xcolor}

\definecolor{colorrojo}{HTML}{C62828}
\definecolor{colorazul}{HTML}{1565C0}
\definecolor{colorverde}{HTML}{2E7D32}
\definecolor{colorgris}{HTML}{757575}
\definecolor{colornegro}{HTML}{222222}

\definecolor{hellgrau}{HTML}{F5F5F5}
\definecolor{rahmen}{HTML}{CCCCCC}
\definecolor{akzent}{HTML}{666666}

% ═══════════════════════════════════════════════════════════════
% BOXEN
% ═══════════════════════════════════════════════════════════════

\usepackage{tcolorbox}
\tcbuselibrary{skins}

\newtcolorbox{lernzielbox}{
    colback=hellgrau, colframe=hellgrau, boxrule=0pt, arc=0pt,
    left=10pt, right=10pt, top=8pt, bottom=8pt,
    before skip=6pt, after skip=8pt
}

\newtcolorbox{grammatikbox}[1][]{
    colback=white, colframe=white, boxrule=0pt, arc=0pt,
    left=10pt, right=6pt, top=6pt, bottom=6pt,
    borderline west={3pt}{0pt}{akzent},
    before skip=4pt, after skip=4pt
}

\newtcolorbox{merkbox}[1][]{
    colback=hellgrau, colframe=hellgrau, boxrule=0pt, arc=0pt,
    left=10pt, right=10pt, top=8pt, bottom=8pt,
    borderline north={2pt}{0pt}{colornegro},
    before skip=6pt, after skip=6pt
}

\newtcolorbox{vokabelbox}[1][]{
    colback=white, colframe=rahmen, boxrule=0.5pt, arc=0pt,
    left=8pt, right=8pt, top=6pt, bottom=6pt,
    before skip=4pt, after skip=4pt
}

\newtcolorbox{tippbox}{
    colback=white, colframe=white, boxrule=0pt, arc=0pt,
    left=8pt, right=4pt, top=3pt, bottom=3pt,
    borderline west={2pt}{0pt}{akzent}
}

% ═══════════════════════════════════════════════════════════════
% BEFEHLE
% ═══════════════════════════════════════════════════════════════

\newcommand{\es}[1]{\seriftext{\textbf{#1}}}
\newcommand{\de}[1]{\textit{\textcolor{akzent}{#1}}}
\newcommand{\gram}[1]{\textsc{#1}}
\newcommand{\abmark}[1]{{\footnotesize\textcolor{akzent}{$\rightarrow$ AB #1}}}
\newcommand{\abschnitt}[2]{\noindent\textbf{\large #1. #2}}
\newcommand{\stamm}[1]{\textcolor{colorrojo}{#1}}

\setlength{\parindent}{0pt}
\setlength{\parskip}{5pt}

% ═══════════════════════════════════════════════════════════════
% KOPF- UND FUSSZEILE
% ═══════════════════════════════════════════════════════════════

\pagestyle{fancy}
\fancyhf{}
\renewcommand{\headrulewidth}{0.5pt}
\renewcommand{\footrulewidth}{0pt}
\lhead{\footnotesize\textsc{Spanisch} · Klasse \klasse}
\rhead{\footnotesize\thema}
\cfoot{\footnotesize\thepage\,/\,\pageref{LastPage}}

% ═══════════════════════════════════════════════════════════════
% DOKUMENT
% ═══════════════════════════════════════════════════════════════

\begin{document}

% --- TITEL ---
\begin{center}
    {\LARGE\bfseries \thema}\\[3pt]
    {\small \untertitel}
\end{center}

\vspace{4pt}

Heute wiederholen wir die Vokabeln und Grammatik aus der letzten Stunde und verbinden sie mit dem \es{pretérito perfecto}. Am Ende kannst du einen Umzugsbericht schreiben.

\vspace{2pt}

\begin{lernzielbox}
\textbf{Das wiederholst und lernst du:}\quad \lernziele
\end{lernzielbox}

% ═══════════════════════════════════════════════════════════════
% ZWEI SPALTEN
% ═══════════════════════════════════════════════════════════════

\begin{multicols}{2}

% ---------------------------------------------------------------
% ABSCHNITT 1: VOKABELWIEDERHOLUNG
% ---------------------------------------------------------------
\abschnitt{1}{Repaso del vocabulario} \de{-- Vokabelwiederholung}

\vspace{4pt}

Erinnerst du dich an die Hausräume und Möbel? Hier ein schneller Test:

\vspace{3pt}

\begin{tippbox}
\footnotesize\textit{Tipp:} Schau auf dein AB von letzter Stunde (Kasten V1), wenn du unsicher bist.
\end{tippbox}

\vspace{3pt}

{{\small Wichtige Räume:}}

\vspace{2pt}

{\footnotesize
\begin{tabular}{@{}ll@{}}
\es{el pasillo} & \de{der Flur} \\
\es{el salón} & \de{das Wohnzimmer} \\
\es{la cocina} & \de{die Küche} \\
\es{el dormitorio} & \de{das Schlafzimmer} \\
\end{tabular}}

\abmark{Aufgabe 1}

\vspace{8pt}

% ---------------------------------------------------------------
% ABSCHNITT 2: BEWEGUNGSVERBEN
% ---------------------------------------------------------------
\abschnitt{2}{Los verbos de movimiento} \de{-- Bewegungsverben}

\vspace{4pt}

Die Bewegungsverben und ihre Präpositionen sind wichtig für den Umzug:

\begin{grammatikbox}
{\small
\begin{tabular}{@{}cll@{}}
$\downarrow$ & \es{bajar de} & \de{herunterbringen von} \\[0.3em]
$\uparrow$ & \es{subir a} & \de{hochbringen nach} \\[0.3em]
$\oplus$ & \es{poner en} & \de{stellen/legen in} \\[0.3em]
$\ominus$ & \es{sacar de} & \de{herausnehmen aus} \\[0.3em]
$\leftarrow$ & \es{traer de} & \de{bringen (her) von} \\[0.3em]
$\rightarrow$ & \es{llevar a} & \de{bringen (hin) nach} \\
\end{tabular}}
\end{grammatikbox}

\begin{merkbox}
\textbf{Merke: Richtungspräpositionen}

\vspace{2pt}
{\small
\textbf{de / del} = von, aus (Richtung weg)\\
\textbf{a / al} = nach, zu (Richtung hin)\\
\textbf{en} = in (Position)}
\end{merkbox}

\abmark{Aufgabe 2}

\columnbreak

% ---------------------------------------------------------------
% ABSCHNITT 3: PRETÉRITO PERFECTO
% ---------------------------------------------------------------
\abschnitt{3}{El pretérito perfecto} \de{-- Perfekt}

\vspace{4pt}

Das \es{pretérito perfecto} beschreibt Handlungen, die bis zur Gegenwart reichen. Du kennst es schon aus Unidad 1!

\begin{grammatikbox}
\textbf{Bildung:} \es{haber} + Partizip

\vspace{4pt}

{\small
\begin{tabular}{@{}ll@{}}
\es{he} & \de{ich habe} \\
\es{has} & \de{du hast} \\
\es{ha} & \de{er/sie hat} \\
\es{hemos} & \de{wir haben} \\
\es{habéis} & \de{ihr habt} \\
\es{han} & \de{sie haben} \\
\end{tabular}}

\vspace{4pt}

{\small
\textbf{Partizip regelmäßig:}\\
-ar → \stamm{-ado}: habl\stamm{ado}\\
-er/-ir → \stamm{-ido}: com\stamm{ido}, viv\stamm{ido}
}

\vspace{4pt}

{\small
\textbf{Partizip unregelmäßig:}\\
poner → \stamm{puesto}\\
hacer → \stamm{hecho}
}
\end{grammatikbox}

\abmark{Kasten G1}

\vspace{8pt}

% ---------------------------------------------------------------
% ABSCHNITT 4: KOMBINATION
% ---------------------------------------------------------------
\abschnitt{4}{Combinación} \de{-- Kombination}

\vspace{4pt}

Jetzt kombinieren wir das pretérito perfecto mit den Bewegungsverben und Objektpronomen:

\vspace{3pt}

{\footnotesize
\es{He puesto la lámpara en el salón.}\\
\de{Ich habe die Lampe ins Wohnzimmer gestellt.}

\vspace{3pt}

\es{La he puesto en el salón.}\\
\de{Ich habe sie ins Wohnzimmer gestellt.}
}

\begin{tippbox}
\footnotesize\textit{Achtung:} Das Objektpronomen steht \textbf{vor} dem konjugierten Verb \es{haber}!
\end{tippbox}

\abmark{Kasten R1 (lesen) + Aufgabe 3}

\end{multicols}

% ═══════════════════════════════════════════════════════════════
% REDEMITTEL
% ═══════════════════════════════════════════════════════════════

\vspace{4pt}

\abschnitt{5}{Comunicación} \de{-- Der Umzugsbericht}

\vspace{4pt}

\begin{vokabelbox}
\begin{multicols}{2}

{\small\textbf{Vergangene Handlungen:}}

\vspace{2pt}

{\footnotesize
\es{He puesto el/la... en...}\\
\de{Ich habe den/die... in... gestellt.}

\vspace{3pt}

\es{Hemos subido los/las... al piso.}\\
\de{Wir haben die... hochgebracht.}

\vspace{3pt}

\es{Mi padre ha sacado...}\\
\de{Mein Vater hat... herausgenommen.}}

\columnbreak

{\small\textbf{Zeitangaben:}}

\vspace{2pt}

{\footnotesize
\es{Esta semana...}\\
\de{Diese Woche...}

\vspace{3pt}

\es{Hoy...}\\
\de{Heute...}

\vspace{3pt}

\es{Esta mañana...}\\
\de{Heute Morgen...}}

\end{multicols}
\end{vokabelbox}

\abmark{Aufgaben 4 + 5}

\vspace{8pt}

% ═══════════════════════════════════════════════════════════════
% ZUSAMMENFASSUNG
% ═══════════════════════════════════════════════════════════════

\begin{merkbox}
\textbf{Zusammenfassung: Was du heute wiederholt und gelernt hast}

\vspace{4pt}

\begin{multicols}{3}
{\footnotesize

\textbf{Wiederholung:}\\
Hausvokabeln\\
6 Bewegungsverben + Präp.\\
4 Objektpronomen

\columnbreak

\textbf{Neu kombiniert:}\\
Pretérito perfecto + Verben\\
Pronomen vor \es{haber}\\
\es{Lo he puesto en...}

\columnbreak

\textbf{Kommunikation:}\\
Umzugsbericht schreiben\\
Vergangenes beschreiben

}
\end{multicols}
\end{merkbox}

\end{document}
